\frenchspacing
\documentclass{amsart}
\usepackage{amssymb}
\usepackage{fancyhdr}
\usepackage{bbm}
\usepackage{enumerate}
\usepackage{graphicx}
\usepackage{tikz}
\pagestyle{fancy}
\allowdisplaybreaks


%% your macros -->
\newcommand{\nn}{\mathbb N}    %% naturals
\newcommand{\zz}{\mathbb Z}    %%integers
\newcommand{\rr}{\mathbb R}    %% real numbers
\newcommand{\cc}{\mathbb C}    %% complex numbers
\newcommand{\ff}{\mathbb F}
\newcommand{\qq}{\mathbb Q}
\newcommand{\cA}{\mathcal{A}}
\newcommand{\cP}{\mathcal{P}}
\newcommand{\cN}{\mathcal{N}}
\newcommand{\cM}{\mathcal{M}}
\newcommand{\cB}{\mathcal{B}}
\newcommand{\cF}{\mathcal{F}}
\newcommand{\cE}{\mathcal{E}}
\newcommand{\cT}{\mathcal{T}}
\newcommand{\limn}{\lim_{n \to \infty}} %%lim n to infty shorthand
\newcommand{\va}{\mathbf{a}}
\newcommand{\vb}{\mathbf{b}}
\newcommand{\vc}{\mathbf{c}}
\newcommand{\vx}{\mathbf{x}}
\newcommand{\vy}{\mathbf{y}}
\newcommand{\vv}{\mathbf{v}}
\newcommand{\vw}{\mathbf{w}}
\newcommand{\vu}{\mathbf{u}}
\DeclareMathOperator{\ee}{\mathbb E}
\DeclareMathOperator{\Var}{Var}  %% variance
\DeclareMathOperator{\Aut}{Aut}
\DeclareMathOperator{\Sym}{Sym}
\DeclareMathOperator{\Cov}{Cov}
\newtheorem{theorem}{Theorem}
\newtheorem{lemma}[theorem]{Lemma}
\newtheorem{corollary}[theorem]{Corollary}
\newtheorem{example}[theorem]{Example}
\theoremstyle{definition}
\newtheorem{definition}{Definition}
\newtheorem{exercise}{Exercise}[section]
\theoremstyle{remark}
\newtheorem{remark}[theorem]{Remark}

\numberwithin{equation}{section}



\begin{document}
\title{Backwards SDEs}
\author{Andrew Lys}
\email{alys@ncsu.edu}

\date{}

\dedicatory{}
\begin{abstract}
    In this paper, we introduce some of the theory of Backwards Stochastic Differential Equations. 
    We begin with the Martingale Representation Theorem, which is a simple example of a BSDE. 
    We then introduce the general theory of BSDEs, and show the connection between BSDEs and PDEs via the nonlinear Feynman-Kac formula. 
    Finally, we discuss some applications of BSDEs in mathematical finance.
\end{abstract}

\maketitle

\section{Introduction}
In this paper, we introduce the theory of Backwards Stochastic Differential Equations (BSDEs). These are extremely useful tools in stochastic analysis, with applications in mathematical finance and stochastic optimal control.
The prototypical BSDE is found in the Martingale Representation Theorem. Akin to how SDEs can be used to generalize stochastic integration to non-linear settings, BSDEs can be used to generalize the Martingale Representation Theorem to non-linear settings.
We will see how the BSDE is extremely natural for pricing and hedging financial derivatives, and how the flexibility and power of BSDEs allows us to price derivatives in more complicated settings than the classical Black-Scholes model.

\subsection{Notation}
Let $(\Omega, \cF_T)$ be a measurable space with a filtration $\cF$ = $(\cF_t)_{t \in [0, T]}$ satisfying the usual conditions. Additionally, let $B_t$ be an $m$-dimensional Brownian motion adapted to $\cF$.
Throughout this paper, we will assume that $\cF = \cF^B$, the natural filtration generated by $B_t$.
We define the following spaces of stochastic processes:
\begin{itemize}
    \item $L^2(\cF, \rr^n)$: the space of $\rr^n$-valued, $\cF$-adapted processes $Y_t$ with finite second moment, i.e. 
    \[
    \ee \int_0^T |Y_t|^2 \mathrm{d}t < \infty
    \]
    \item 
    \(
    L^{p, q} (\cF, \rr^n) = \{Z \in L^0(\cF, \rr^n) : \|Z\|_{L^p([0,T])} \in L^q(\cF_T)\}
    \)
\end{itemize}
\section{Martingale Representation Theorem}
An example of a simple BSDE is in the form of the martingale representation theorem. The technique for solving it is very similar to how we solve backwards PDEs. 
\begin{theorem}
    Let $\xi$ be a $\cF_T$-measurable random variable with finite second moment. Then there exists a unique adapted process $\sigma \in L^2(\cF, \rr^d)$ such that 
    \begin{equation}
        \xi = \ee \xi + \int_0^T \sigma_t \cdot \mathrm{d}B_t
    \end{equation}
\end{theorem}
This is exactly a backwards SDE. The proof will follow in 5 parts, of which we will only prove the first step, since the rest are basic approximations. 
\begin{proof}
    \begin{enumerate}[i)]
        \item Suppose $\xi = g(B_T)$ for $g \in C^2_b$.
        \item Suppose $g$ is borel measurable and bounded and that $\xi = g(B_T)$.
        \item Suppose $g$ is borel measurable and bounded and there exist $0 \le t_1 < \ldots < t_n \le T$ such that $\xi = g(B_{t_1}, \ldots, B_{t_n})$.
        \item Suppose $\xi$ is essentially bounded and $\cF_T^B$-measurable.
        \item Suppose $\xi$ is $\cF_T^B$-measurable with finite second moment.
    \end{enumerate}
    We will prove part (i). Recall the semigroup operator corresponding to Brownian motion defined by
    \begin{equation}
        (S_t f)(x) = \ee_x f(B_t)
    \end{equation}
    We are looking for a process $M_t$ such that $M_T = \xi$ and $\mathrm{d}M_t = \sigma_t \mathrm{d}B_t$. A very place to start looking for $M_t$ is, as we are used to doing with backwards PDEs, is to let $u$ be given by
    \begin{equation}
        u(t, x) = (S_{T-t} g)(x)
    \end{equation}
    Note that $u$ satisfies the backwards heat equation. By Kolmogorov's backwards equation, we have that
    \begin{align*}
        \partial_t u &= - \partial_s S_{s} g |_{s = T-t} = - \frac{1}{2} \Delta_x S_{T-t} g \\
        &= - \frac{1}{2} \Delta_x u
    \end{align*}
    Plugging in $t = T$ gives us:
    \[
    u(T, x) = g(x)
    \]
    Which, when combined with $x = B_T$ gives us 
    \[
    u(T, B_T) = g(B_T) = \xi
    \]
    Naturally, we define $M_t = u(t, B_t)$. Applying It\^o's formula to $M_t$ gives us
    \begin{align*}
        \mathrm{d} M_t &= \partial_t u \mathrm{d} t + \nabla_x u \cdot \mathrm{d}B_t + \frac{1}{2} \Delta_x u \mathrm{d}t \\
        &= \nabla_x u \cdot \mathrm{d}B_t
    \end{align*}
    Which gives us $\sigma$ for free!
    \begin{equation}
        \sigma_t = \nabla_x u(t, B_t) = \nabla_x \ee_x (g(B_{T-t})) |_{x = B_t}
    \end{equation}
    Since $g \in C^2_b$, we have that $\sigma$ is adapted and in $L^2(\cF, \rr^d)$. This completes part (i).
    The rest of the parts follow from uninteresting but basic approximation arguments.
\end{proof}
\begin{remark}
    This $\sigma_t$ that we found has an important interpretation in mathematical finance. If $M_t$ is the value of a contingent claim at time $t$, then $\sigma_t = \nabla_x M_t$ is the amount of each asset we need to hold at time $t$ in order to perfectly hedge the claim.
\end{remark}
We have the following simple corollary which is the more commonly seen version of the martingale representation theorem.
\begin{corollary}
    Let $M_t$ be a martingale with respect to $\cF$ with finite second moment. Then, there exists a unique $\cF$-adapted process $\sigma$ with finite second moment such that
    \begin{equation}
        M_t = \ee M_0 + \int_0^t \sigma_s \cdot \mathrm{d}B_s
    \end{equation}
\end{corollary}
\begin{proof}
    By Theorem 1, we get $\sigma$ such that 
    \[
    M_T = \ee M_T + \int_0^T \sigma_s \cdot \mathrm{d}B_s
    \]
    Then, let 
    \[
    \tilde{M}_t := \ee M_T + \int_0^t \sigma_s \cdot \mathrm{d}B_s
    \]
    Then by the properties of the It\^o integral, $\tilde{M}_t$ is a martingale with respect to $\ff^B$ with finite second moment. 
    Then we have
    \[
    M_t = \ee[M_T | \cF_t^B] = \ee[\tilde{M}_T| \cF_t^B] = \tilde{M}_t
    \]
    giving us $M_0 = \tilde{M}_0 = \ee M_T$ and completing the proof.
\end{proof}

\section{Backward Stochastic Differential Equations}

The introduction of BSDEs is conceptually very similar to the introduction of SDEs. 
It is very easy to solve a linear SDE with just stochastic integration. 
This is essentially the role the Martingale Representation Theorem plays for BSDEs. 
However, when we introduce a non-linear term, we require the concept of an It\^o process to solve our problem. 
Similarly, BSDEs deal with the non-linear case. 

For an example of the Martingale Representation Theorem as a BSDE, we have the following example.
\begin{example}
    Let $\xi \in L^2(\cF_T)$. Then we have $Y_t = \ee[\xi | \cF_t]$ is a martingale. By the Martingale Representation Theorem, there exists $Z \in L^2(\cF)$ such that 
    \[
    \begin{cases}
        \mathrm{d}Y_t = Z_t \cdot \mathrm{d} B_t \qquad t \in [0, T] \\
        Y_T = \xi
    \end{cases}
    \]
\end{example}
We always solve BSDEs with a pair of processes $(Y, Z)$. However, in some idea, $Z$ is determined by $Y$, and can be conceptually thought of as the derivative of $Y$ with respect to the Brownian motion.
As we showed above, Backwards SDEs can be solved in a similar manner to backwards PDEs. 
However, the difficulty of BSDEs comes in verifying the adaptedness of solution processes. 
This is also why $Z_t$ is so important. For instance, suppose $Z_t \equiv 0$. Then $Y_t = \xi$ for all $t$. 
Unless $\xi$ is $\cF_0^B$-measurable, $Y_t$ is not adapted, and is not a solution we want.  

We now consider the various kinds of BSDEs. First, we propose the general nonlinear BSDE. 
\begin{definition}
    Let $Y \in L^2(\cF, \rr^n)$, $Z \in L^2(\cF, \rr^{n \times m})$, $B_t$ be an $m$-dimensional Brownian motion, adapted to $\cF$, $\xi \in L^2(\cF_T, \rr^n)$, and $f: [0, T] \times \Omega \times \rr^n \times \rr^{n \times m} \to \rr^n$ be measurable and adapted. 

    A \textbf{backward stochastic differential equation} (BSDE) is an equation of the form
    \begin{equation}
        Y_t = \xi + \int_t^T f_s(Y_s, Z_s) \mathrm{d}s - \int_t^T Z_s \cdot \mathrm{d}B_s \qquad t \in [0, T]
    \end{equation}
\end{definition}

Additionally, for well-posedness in basic situations we usually additionally assume the following:
\begin{enumerate}[i)]
    \item $f$ is uniformly Lipschitz in $y$ and $z$. 
    \item $f^0 := f(0, 0) \in L^{1,2}(\cF, \rr^n)$.
\end{enumerate}
In fact, we have the following well-posedness result for BSDEs.
\begin{theorem}
    Under the assumptions above, there exists a unique solution $(Y, Z) \in L^2(\cF, \rr^n) \times L^2(\cF, \rr^{n \times m})$ to the BSDE in Definition 1.
\end{theorem}
Now, we look at a simple case of a BSDE, the linear case. First, we look at the case where $f$ is linear.
\begin{theorem}
    The following BSDE has a unique solution $(Y, Z) \in L^2(\cF, \rr^n) \times L^2(\cF, \rr^{n \times m})$, and $Y$ is continuous almost surely. 
    \begin{equation}
        Y_t = \xi + \int_t^T f_s^0 \mathrm{d} s - \int_t^T Z_s \cdot \mathrm{d}B_s
    \end{equation}
\end{theorem}
\begin{proof}
    This follows almost directly from the Martingale Representation Theorem. 
    Our candidate for $Y_t$ is very naturally going to be 
    \begin{equation}
        Y_t := \ee \left[
            \xi + \int_t^T f_s^0 \mathrm{d} s \bigg| \cF_t
        \right]
    \end{equation}
    Note that 
    \[
    \tilde{Y}_t := \ee \left[
        \xi + \int_0^T f_s^0 \mathrm{d} s \bigg| \cF_t
    \right]
    \]
    is a square integrable martingale as we can easily see from the following:
    \begin{align*}
        \int_0^T \ee |\tilde{Y}_t|^2 \mathrm{d} t & = \int_0^T \ee \left| \ee \left[
            \xi + \int_0^T f_s^0 \mathrm{d} s \bigg| \cF_t \right] \right|^2 \mathrm{d} t \\
            &\le \int_0^T \ee \ee \left[
                \left|
                    \xi + \int_0^T f_s^0 \mathrm{d}s 
                \right|^2 \bigg| \cF_t
            \right]\\
            &\le 2T \left(
                \ee |\xi|^2 + \ee \left|\int_0^T f_s^0 \mathrm{d} s \right|^2 
            \right) \\
            &\le 2T \ee |\xi|^2 + 2T^2 \ee \int_0^T |f_s^0|^2 \mathrm{d} s < \infty
    \end{align*}
    Therefore, $\tilde{Y}_t$ is a square integrable martingale. By the MRT, we have
    \[
    \mathrm{d} \tilde{Y}_t = Z_t \cdot \mathrm{d}B_t
    \]
    And note that we have:
    \[
    \tilde{Y}_t = \ee \left[
        \xi + \int_0^t f_s^0 \mathrm{d} s + \int_t^T f_s^0 \mathrm{d} s \bigg| \cF_t
    \right] = \int_0^t f_s^0 \mathrm{d} s + Y_t
    \]
    Which gives us
    \[
    \tilde{Y}_T - \tilde{Y}_t = Y_T - Y_t + \int_t^T f_s^0 \mathrm{d} s = \int_t^T Z_s \cdot \mathrm{d}B_s
    \]
    Rearranging gives us the desired BSDE:
    \[
    Y_t = \xi + \int_t^T f_s^0 \mathrm{d} s - \int_t^T Z_s \cdot \mathrm{d}B_s
    \]
\end{proof}

We now consider the more general linear case. We propose the exact solution without proof, as the proof is very similar to the case in Theorem 5. 
\begin{theorem}
    Let $\alpha \in L^\infty (\cF, \rr)$, $\beta \in L^\infty(\cF, \rr^{m})$ and $f^0 \in L^{1,2}(\cF, \rr^n)$. Then the following BSDE
    \begin{equation}
        Y_t = \xi + \int_t^T \alpha_s Y_s + Z_s \beta_s + f_s^0 \mathrm{d} s - \int_t^T Z_s \cdot \mathrm{d}B_s
    \end{equation}
    has a unique solution $(Y, Z) \in L^2(\cF, \rr^n) \times L^2(\cF, \rr^{n \times m})$, and $Y$ is continuous almost surely. Furthermore, the solution is given by
    \begin{equation}
        Y_t = \Gamma_t^{-1} \ee \left[
            \Gamma_T \xi + \int_t^T \Gamma_s f_s^0 \mathrm{d} s \bigg| \cF_t
        \right]
    \end{equation}
    Where $\Gamma_t$ is given by
    \[
    \Gamma_t = \exp\left(\int_0^t \alpha_s \mathrm{d} s\right) \cE\left(
        \int_0^t \beta_s \cdot \mathrm{d} B_s
    \right)
    \]
\end{theorem}
Most nonlinear BSDEs do not have explicit solutions, and require numerical methods and well-posedness results. 

\section{Nonlinear Feynman-Kac Formula}
We now look at the connection between BSDEs and PDEs. Recall Dynkin's formula and the Feynman-Kac formula for SDEs. 
Namely, if we have the second order elliptic differential operator $A$ given by
\[
Af = b \cdot \nabla f + \frac12 \sigma\sigma^\intercal \cdot \nabla^2 f
\]
Then the solution to the following parabolic PDE
\begin{equation}
    \begin{cases}
        Au - \partial_t u = 0 & t > 0\\
        u(0, x) = f(x)  & x\in \rr^n
    \end{cases}
\end{equation}
is given by Dynkin's formula, where we have the following SDE:
\begin{equation}
\mathrm{d} X_t = b(X_t) \mathrm{d} t + \sigma(X_t) \mathrm{d} B_t
\end{equation}
Then the solution to the PDE is given by
\begin{equation}
    u(t, x) = \ee_x f(X_t)
\end{equation}
Additionally, if we have the following parabolic PDE with a potential term, 
\begin{equation}
    \begin{cases}
        Av - \partial_t v - q v = 0 & t > 0 \\
        v(0, x) = f(x) & x \in \rr^n 
    \end{cases}
\end{equation}
Then the solution is given by the Feynman-Kac formula:
\begin{equation}
    v(t, x) = \ee_x \left[
        f(X_t) \exp\left(
            -\int_0^t q(X_s) \mathrm{d} s
        \right)
    \right]
\end{equation}

These are very powerful results in the field of stochastic processes, as they allow us to solve certain PDEs usign stochastic methods. 
However, they are quite limited in scope, as they only deal with nice linear PDEs, and do not readily extend to nonlinear PDEs. 

As Feynman-Kac can be applied to nice linear PDEs, BSDEs can be applied to not-so-nice nonlinear PDEs. 
First, we introduce the following decoupled Forward-Backward SDE (FBSDE) with deterministic coefficients.
\begin{equation}
    \begin{cases}
        X_t = x + \int_0^t b(s, X_s) \mathrm{d} s + \int_0^t \sigma(s, X_s) \mathrm{d} B_s \\
        Y_t = g(X_T) + \int_t^T f(s, X_s, Y_s, Z_s) \mathrm{d} s - \int_t^T Z_s \cdot \mathrm{d} B_s
    \end{cases}
\end{equation}
Under the following regularity assumptions, this problem is well-posed. 
\begin{enumerate}[i)]
    \item $b, \sigma, f, g$ are deterministic, and $b(\cdot, 0), \sigma(\cdot, 0), f(\cdot, 0, 0, 0) , g(0)$ are bounded.
    \item $b, \sigma, f, g$ are uniformly Lipschitz in all their spatial variables.
\end{enumerate}
The following additional assumption will ensure stability in numerical schemes.
\begin{enumerate}
    \item[iii)] $b, \sigma, f$ are uniformly H\"older-$\frac12$ continuous in $t$.
\end{enumerate}
Skipping the proofs, we have the following nonlinear Feynman-Kac formula.
\begin{theorem}
    Assume the above assumptions hold and $u \in C^{1, 2}([0, T] \times \rr^n, \rr^m)$ is a classical solution to the following PDE:
    \begin{equation}
        \begin{cases}
            \partial_t u + A u + f(t, x, u, \nabla_x u\sigma) = 0\\
            u(T, x) = g(x)
        \end{cases}
    \end{equation}
    Then 
    \begin{equation}
        Y_t = u(t, X_t), \quad Z_t = \nabla_x u(t, X_t) \sigma(t, X_t)
    \end{equation}
\end{theorem}


\section{Applications}
BSDEs have many applications in finance and stochastic optimal control. 
We will briefly discuss some of the financial applications. 

As we have seen above a BSDE is solved with a pair of processes $(Y, Z)$.
In financial applications, $(Y, Z) = (V, \Delta)$ where $V$ is the value of a portfolio and $\Delta = h \sigma S$ is the hedge. 
More general BSDEs invoke a nonlinear running cost function $g(t, Y, Z)$. 
If $g$ is linear, we essentially invoke the martingale representation theorem as we have used for the linear BSDE above. 
However, if $g$ is nonlinear, as is often the case in finance (spread between borrowing and lending rates), we require Picard iteration or some sort of contraction mapping argument. 

Nonlinearity can appear not just from $g$ but also from early exercise features present in American options, or two way early exercise features present in game options.
In this case, we require a reflected BSDE or doubly reflected BSDE. 

\subsection{Optimal Hedging and Replication}
Let $\xi$ be a contingent claim maturing at time $T$. Then we may price the claim by finding a pair of processes $(\lambda_t, h_t)$ that replicate the payoff $\xi$. 
$\lambda_t$ is the amount invested in a risk-free asset, and $h_t$ is the amount invested in the one risky asset. 
We assume that the risk-free rate is $r > - 1$.
In order to correctly price the claim, we require that our portfolio is self-financing, meaning we cannot add or withdraw money from the portfolio. 
This gives us the following SDE for the value of the portfolio $V_t$:
\begin{equation}
    \mathrm{d} V_t = r \lambda_t \mathrm{d} t + h_t \mathrm{d} S_t
\end{equation}
If we further assume that the risky assets are driven by a geometric Brownian motion with drift $\mu$ and volatility $\sigma$, we have
\begin{align}
    \mathrm{d} V_t &= \left[r \lambda_t + h_t  S_t \mu \right]\mathrm{d} t  + h_t  S_t \sigma \mathrm{d} B_t 
\end{align}
Now recall that the value of our portfolio is given by 
\begin{equation}
    V_t = \lambda_t e^{r t} + h_t S_t
\end{equation}
Which gives us the following SDE for $V_t$:
\begin{equation}
    \mathrm{d} V_t = \left[
        r (V_t - h_t S_t) + h_t S_t \mu 
    \right] \mathrm{d} t + h_t S_t \sigma \mathrm{d} B_t
\end{equation}
Writing $Y_t := V_t$ and $Z_t := \sigma S_t h_t$, we have $Y_T = \xi$ and we exactly get a linear BSDE for the value of the portfolio!
\begin{equation}
    \begin{cases}
        \mathrm{d} Y_t = \left[r \left[
             Y_t - \frac{Z_t}{\sigma S_t}
        \right] + \frac{\mu Z_t}{\sigma S_t}\right]\mathrm{d} t + Z_t \mathrm{d} B_t \qquad t \in [0, T] \\
        Y_T = \xi
    \end{cases}
\end{equation}
Since this is a linear BSDE, we can explicitly solve for the value of the portfolio $Y_t$ using Theorem 6. 
In this situation we have that $Y_0$ is the arbitrage-free price of the contingent claim $\xi$ at time $0$, and $h_t = \frac{Z_t}{\sigma S_t}$ is the amount of the risky asset we need to hold at time $t$ in order to perfectly hedge the claim.
We then automatically get $\lambda_t = V_t - \frac{Z_t}{\sigma}$

This is a very simple example of how BSDEs can be used in mathematical finance. In fact, in this case BSDE theory is overkill, as this solution can be derived from the simple Black-Scholes model. 
However, with even modest extensions to this model, we already require BSDE theory to solve the problem. 

\begin{example}
    Consider the same setup as above, but now suppose that the lending and borrowing rates are different, with $r_l < r_b$. 
    Our self-financing condition now becomes
    \begin{equation}
        \mathrm{d} V_t = [
            r_l (V_t - h_t S_t)^+ - r_b (V_t - h_t S_t)^- 
        ] \mathrm{d} t + h_t \mathrm{d} S_t
    \end{equation}
    This BSDE is now nonlinear.
\end{example}

\subsection{American Options}
We now consider a more complicated financial derivative that has an early exercise feature, an American option.

Our model is similar to the previous section, with a risk-free asset with and a risky asset driven by a geometric Brownian motion.
To simplify things, we will consider the case of $\mu = r = 0$. Under these assumptions the price of a European option is 
\[
Y_0 = \ee [\xi]
\]
We consider the American option, with payoff process $L_t$. For example, the payoff process for an American call option is $L_t = (S_t - K)^+$

The holder of the American option can choose to exercise the option at any stopping time $\tau \in \cT$, where $\cT$ is the set of all stopping times taking values in $[0, T]$.
Therefore, the price of the American option at time $0$ is given by
\begin{equation}
    Y_0 = \sup_{\tau \in \cT} \ee [L_\tau]
\end{equation}
Let 
\begin{equation}
    Y_t := \mathrm{ess}\sup_{\tau \in \cT^t} \ee [L_\tau | \cF_t]
\end{equation}
Where 
\begin{equation}
    \cT^t := \{\tau \in \cT : t \le \tau \le T\}
\end{equation}
The process $Y_t$ is called the Snell envelope of $L_t$. It is important to use the essential supremum here, as the supremum might involve uncountably many stopping times.
The Snell envelope has the following important properties:
\begin{enumerate}[i)]
    \item $Y_t \ge L_t$ almost surely for all $t \in [0, T]$. This is because $\tau = t \in \cT$.
    \item $Y_t$ is a supermartingale. If we were dealing with a European option, the the price process would be a martingale. However, an American option has a decreasing non-negative time value, so the price process is a supermartingale.
\end{enumerate}
The second property is very important, as it allows us to use the Doob-Meyer decomposition theorem.
By the Doob-Meyer decomposition theorem, we have that there exists a unique decomposition of $Y_t$ into a martingale $M_t$ and a predictable, increasing process $K_t$ with $M_0 = K_0 = 0$ such that
\[
Y_t = Y_0 + M_t - K_t
\]
By the Martingale Representation Theorem, there exists $Z_t$ such that
\[
\mathrm{d} M_t = Z_t \mathrm{d} B_t
\]
By combining these two results, we have the following BSDE for the price process of the American option:
\begin{align*}
    Y_T &= L_T = Y_0 + M_T - K_T = Y_0 + M_t + \int_t^T Z_s \mathrm{d}B_s - K_T\\
    &= Y_t + \int_t^T Z_s \mathrm{d} B_s - (K_T - K_t) \\
    \implies Y_t &= L_T + \int_t^T \mathrm{d} K_s - \int_t^T Z_s \mathrm{d} B_s
\end{align*}

Notice that when $Y_t > L_t$, there would be an arbitrage opportunity if one were to exercise the option. 
Therefore, when $Y_t > L_t$, we must have $\mathrm{d} K_t = 0$, which gives us the following condition, called the Skorokhod condition:
\begin{equation}
    \int_0^T (Y_t - L_t) \mathrm{d} K_t = 0
\end{equation}

Putting this all together, we have what is known as a reflected BSDE for the price process of the American option:
\begin{equation}
    \begin{cases}
        Y_t = L_T + \int_t^T \mathrm{d} K_s - \int_t^T Z_s \mathrm{d} B_s \\
        Y_t \ge L_t \\
        \int_0^T (Y_t - L_t) \mathrm{d} K_t = 0
    \end{cases}
\end{equation}
This is termed as a reflected BDSE because of the lower barrier that $L_t$ provides for $Y_t$. This is defined in analogy to the reflection principle for random walks and Brownian motion.

We define the general reflected BSDE as follows. 
\begin{equation}
    \begin{cases}
        Y_t = \xi + \int_t^T f_s(Y_s, Z_s) \mathrm{d} s + \int_t^T \mathrm{d} K_s - \int_t^T Z_s \cdot \mathrm{d} B_s \\
        Y_t \ge L_t \\
        \int_0^T (Y_t - L_t) \mathrm{d} K_t = 0
    \end{cases}
\end{equation}
A solution to an RBSDE is a triple $(Y, Z, K)$ where $Y$ is continuous and $K$ is increasing and continuous. 

If we consider an option with both a lower and upper barrier, we get a doubly reflected BSDE. This arises in game options, where both the buyer and seller have the option to exercise early.
\bibliographystyle{amsplain}
\bibliography{refs}


\end{document}